\BeleskeID{04-s026}
% element: 04-s026:e01 — Pitanje binarne tačke i mogućnost celobrojnog, razlomljenog ili mešovitog zapisa
% element: 04-s026:e02 — Registar 01101000 i nemogućnost posebnog koda tačke sa 0 i 1
% element: 04-s026:e03 — Projektant predviđa položaj; oba podatka o značenju moraju biti poznata
% element: 04-s026:e04 — Primeri tačke iza D5, iza D0, ispred D7 i drugi položaji
% element: 04-s026:d01

Binarna tačka u prikazu sa fiksnom tačkom nije dodatna fizička ćelija registra. Njeno mesto je deo dogovora o tumačenju svih bitova. Zato ista reč može predstavljati različite vrednosti, iako je njen fizički sadržaj potpuno isti.

Ako je tačka iza $D_5$, pet bitova $D_4$ do $D_0$ nalazi se u razlomljenom delu. Ako je iza $D_0$, nema razlomljenih bitova. Ako je ispred $D_7$, svih osam mesta ima razlomljene težine. Potrebno je znati i da sadržaj predstavlja broj i koji je od tih položaja dogovoren.
